\documentclass[11pt]{article}

%% ── packages ──────────────────────────────────────────────────────
\usepackage[margin=1in]{geometry}
\usepackage{amsmath,amssymb,amsthm}
\usepackage{mathtools}
\usepackage{stmaryrd}
\usepackage{mathrsfs}
\usepackage{enumitem}
\usepackage{booktabs}
\usepackage{array}
\usepackage{xcolor}
\usepackage[protrusion=true,expansion=false]{microtype}
\usepackage[colorlinks=true,linkcolor=NavyBlue,citecolor=NavyBlue,urlcolor=SkyBlue]{hyperref}

%% ── F1R3FLY brand colours ─────────────────────────────────────────
\definecolor{NavyBlue}{HTML}{1A1A2E}
\definecolor{SkyBlue}{HTML}{3FA9F5}
\definecolor{Yellow}{HTML}{F3D630}
\definecolor{Sage}{HTML}{8BB999}

%% ── theorem environments ──────────────────────────────────────────
\theoremstyle{definition}
\newtheorem{definition}{Definition}[section]
\newtheorem{construction}[definition]{Construction}
\newtheorem{example}[definition]{Example}
\newtheorem{remark}[definition]{Remark}
\newtheorem{notation}[definition]{Notation}

\theoremstyle{plain}
\newtheorem{theorem}[definition]{Theorem}
\newtheorem{proposition}[definition]{Proposition}
\newtheorem{lemma}[definition]{Lemma}
\newtheorem{corollary}[definition]{Corollary}
\newtheorem{axiomx}[definition]{Property}

%% ── notation macros (shared with the Cost note) ──────────────────
\newcommand{\GSLT}{\mathrm{GSLT}}
\newcommand{\iGSLT}{\mathbf{iGSLT}}
\newcommand{\ciGSLT}{\mathbf{ciGSLT}}
\newcommand{\CostC}{\mathbf{Cost}\text{-}\ciGSLT}
\newcommand{\CostCp}{\mathbf{Cost}\text{-}\ciGSLT_{\bullet}}
\newcommand{\Cost}{\mathfrak{C}}
\newcommand{\near}{\mathsf{near}}
\newcommand{\compute}{\mathsf{compute}}
\newcommand{\Kp}{K_{\mathrm{p}}}
\newcommand{\Ke}{K_{\mathrm{e}}}
\newcommand{\Kctx}[1]{K([\,],\, #1)}
\newcommand{\lstk}[2]{S(#1,\, #2)}     % located stack S(I, content)
\newcommand{\cons}{\mathbin{::}}        % temporal cons
\newcommand{\Seq}{\equiv}
\newcommand{\red}{\longrightarrow}
\newcommand{\Sig}{\mathsf{Sig}}
\newcommand{\hashf}{\#}
\newcommand{\cf}{\mathsf{cf}}
\newcommand{\Term}{\mathcal{T}}
\newcommand{\id}{\mathrm{id}}

%% ── notation macros (new to this note) ────────────────────────────
\newcommand{\foot}{\partial}                       % footprint = write-set (consumed)
\newcommand{\rfoot}{\partial^{\mathsf r}}          % read-set (depended on, not consumed)
\newcommand{\fire}[1]{\xrightarrow{\;#1\;}}
\newcommand{\Runs}{\mathrm{Runs}}
\newcommand{\Hist}{\mathrm{Hist}}
\newcommand{\Conf}{\mathrm{Conf}}
\newcommand{\Ev}{\mathbb{E}}
\newcommand{\Evt}{E}
\newcommand{\dpast}[1]{{\downarrow}#1}
\newcommand{\lc}{\lessdot}                          % covering / link
\newcommand{\MeasCaus}{\mathbf{MeasCaus}}
\newcommand{\Caus}{\mathbf{Caus}}
\newcommand{\Map}{\mathcal{M}}
\newcommand{\co}{\mathbin{\mathsf{co}}}            % concurrency
\newcommand{\confl}{\mathrel{\#}}                  % conflict
\newcommand{\reach}{\mathsf{Reach}}
\newcommand{\linz}{\mathsf{lin}}
\newcommand{\spent}{\mathsf{phlo}}
\newcommand{\meas}{\mu}
\newcommand{\propertime}{\tau}
\newcommand{\restrict}{\!\restriction\!}
\newcommand{\defeq}{\;\mathrel{\vcentcolon=}\;}
\newcommand{\Surf}{\mathit{Surf}}

\theoremstyle{remark}

\title{\textbf{Spacetime from Cost}\\[2pt]
\large A functor from cost-accounted ciGSLTs to measured causal sets\\[6pt]
\normalsize\textit{Companion note to ``Continued Interactive GSLTs and the Cost Endofunctor''}}

\author{L.~G.~Meredith\thanks{\texttt{lgreg.meredith@f1r3fly.io}. This
note was developed in dialogue with Claude (Anthropic); the framing of
located authority as the Lorentzian welding of $K$ and $\cons$, and the
identification of the spacetime invariant with hereditary
history-preserving bisimulation, emerged in that exchange.}\\[2pt]
\small F1R3FLY.io}

\date{June 2026}

\begin{document}
\maketitle

\begin{abstract}
The cost endofunctor $\Cost$ of~\cite{CostNote} equips an interactive
graph-structured lambda theory with two combination operators of
opposite character: a spatial contact constructor $K$, recording what is
near what, and a temporal cons $\cons$ on located resource stacks,
recording what is spent next. That note observed, in its closing
correspondences, that when $K$ carries equations the surface at which a
cut fires must be carried nominally by a nearness operator, and that the
resulting \emph{located} stack $\lstk{I}{-}$ ties a temporal resource to
a spatial surface, so that authority entangles space and time into ``a
spacetime.'' Here we make that map precise. We define a functor $\Map$
from pointed cost-accounted continued interactive GSLTs to
\emph{measured causal sets} --- locally finite partial orders with the
counting measure, the objects that causal-set quantum gravity takes as
discrete spacetimes via Sorkin's slogan \emph{Order $+$ Number $=$
Geometry}~\cite{BLMS,SorkinValdivia}. Events are cut firings; the causal
order is enabling; concurrency is Bernstein independence of footprints,
while all nondeterminism (races) is carried by conflict; the measure is
spent phlogiston. We prove: that every history of $\Cost(G)$ is a causal
set (local finiteness); that the counting measure of an event's causal
past equals the phlogiston spent to reach it (Number $=$ Volume), via the
modulus proposition of~\cite{CostNote}; that discrete proper time is the
longest chain and satisfies the Lorentzian reverse triangle inequality;
that a single located stack contributes a timelike chain while non-near
surfaces contribute spacelike antichains, so that located authority is
exactly what forces a $(1{+}n)$ ordering fraction rather than a
degenerate chain or antichain; that the construction is functorial and
that the measured causal set is an invariant of \emph{hereditary
history-preserving} bisimulation (and, pointedly, \emph{not} of ordinary
interleaving bisimulation, which forgets the concurrency that is the
space); and that one reduction step is one cover in the configuration
domain whose new links are a local functional of the consumed footprint
--- the discrete shape of the field-equation coupling. The continuum
Lorentzian manifold appears only as the Hauptvermutung class of a
manifoldlike causal set, the fourth and idealizing quotient. Throughout
we separate what is proved here from the standard external theorems we
invoke.
\end{abstract}

\tableofcontents
\bigskip

%% ══════════════════════════════════════════════════════════════════
\section{Introduction}
\label{sec:intro}

The note~\cite{CostNote} closed with two correspondences to physical
theories, offered in the spirit of Curry--Howard as resemblances
detailed enough to invite investigation. The first read a $K$-tree as a
distribution of information, cut elimination as a dynamics that changes
adjacency, and the coupling between the two --- the information
distribution shaping the nearness structure, the nearness structure
selecting the next redistribution --- as the shape of Einstein's
intuition. It then observed that object capabilities sharpen rather than
dissolve the picture: a \emph{located} resource stack $\lstk{I}{-}$ ties
a temporal resource (what may be spent next) to a spatial surface (who is
near), so a draw requires both spatial proximity ($\near(I,J)$ defined)
and temporal availability (a token at the head of the stack). Space and
time can then no longer be varied independently; authority
\emph{entangles} them, and ``what was a product of a space and a time
becomes a single coupled structure gating interaction --- a spacetime.''

That note left the formal development to future work. This note supplies
it, in the one direction that is both precise and already standard: the
target is not a manifold but a \emph{causal set}. The causal-set
programme of Bombelli, Lee, Meyer, and Sorkin~\cite{BLMS} takes a locally
finite partial order, equipped with the counting measure, as a discrete
spacetime in its own right, recovering the continuum only as an
approximating limit. Its founding slogan, \emph{Order $+$ Number $=$
Geometry}~\cite{SorkinValdivia}, is exactly what the two monoids of
$\Cost$ supply: the causal order is the enabling order of cut firings,
and the number --- the volume --- is the spent phlogiston. A Lorentzian
manifold, when it exists, is then the Hauptvermutung class of the causal
set: the family of manifolds into which it faithfully embeds, unique up
to approximate isometry. This is the precise content of the user's
``(quotiented) spacetime manifold,'' and it factors the
\emph{quotiented} into four named identifications
(Section~\ref{sec:quotients}).

\paragraph{What is proved and what is cited.}
We are explicit about the boundary. Proved here, from the operational
semantics of $\Cost(G)$ and the structural results of~\cite{CostNote},
are: the diamond/locality of independent firings
(Lemma~\ref{lem:diamond}); the prime-event-structure axioms
(Lemma~\ref{lem:pes}); local finiteness, i.e.\ that every history is a
causal set (Theorem~\ref{thm:causet}); Number $=$ Volume $=$
phlogiston (Theorem~\ref{thm:numbervolume}); the Lorentzian reverse
triangle inequality for proper time (Proposition~\ref{prop:revtriangle});
the signature criterion relating located authority to the ordering
fraction (Proposition~\ref{prop:signature}); functoriality and the
HHP-bisimulation invariance of the spacetime
(Theorem~\ref{thm:functor}, Corollary~\ref{cor:invariance}); and the
naturality of the cover along reduction
(Proposition~\ref{prop:naturality}). Cited as standard external results,
not reproved, are: the representation of asynchronous transition systems
and trace languages by prime event structures and prime algebraic
domains~\cite{NPW,WinskelES,WinskelNielsen,Mazurkiewicz,Bednarczyk}; the
characterisation of event-structure isomorphism by hereditary
history-preserving bisimulation~\cite{vGG}; the Myrheim--Meyer dimension
estimator~\cite{Myrheim,MeyerThesis}; and the closeness/Hauptvermutung
theory of causal sets and Lorentzian
manifolds~\cite{BombelliNoldus,MajorRideoutSurya}. We use the
external results only to package and to interpret; every claim that
depends on the specifics of $\Cost$ is proved.

%% ══════════════════════════════════════════════════════════════════
\section{Recollection: the cost-accounted theory}
\label{sec:recall}

We recall just enough of~\cite{CostNote} to be self-contained, fixing the
notation this note operates on. Let $G$ be a continued interactive GSLT:
an interactive theory presented as an \emph{interaction cut}. Its
dynamics brings two introductions into contact along matching interaction
surfaces and contracts them,
\begin{equation}
\label{eq:cut}
  K\bigl(\Kp(I, C_{\mathrm p}),\; \Ke(J, C_{\mathrm e})\bigr)
  \;\red\;
  \compute(I, J, C_{\mathrm p}, C_{\mathrm e}),
\end{equation}
where $K$ is the contact constructor, $\Kp(I,C_{\mathrm p})$ and
$\Ke(J,C_{\mathrm e})$ the program- and environment-side introductions,
$I,J$ their surfaces, and $\compute$ Milner's pseudo-application.

The cost endofunctor $\Cost$ produces from $G$ a metered calculus
$\Cost(G)$: it signs terms with elements of a commutative signature
monoid $(\Sig,*,()\,)$ via $\hashf=\mathrm{digest}\circ\cf$, adjoins a
\emph{temporal token stack}, makes $K$ polymorphic, and replaces the cut
by a \emph{gated} family of rules. The temporal stack is a free monoid
(a list) under cons $\cons$; it is never commutative, because reduction
is sequential. Following Section~11 of~\cite{CostNote}, when $K$ carries
any equation the stack is \emph{located} at the surface that may draw on
it, written
\[
  \lstk{I}{s \cons S'},
\]
a stack indexed by the surface $I$, with head token $s$ and tail $S'$.
The metered reduction is the gated cut: schematically, writing
$\mathcal{C}[-]$ for a configuration context,
\begin{equation}
\label{eq:gatedcut}
  \mathcal{C}\Bigl[\,
     K\bigl(\Kp(I, C_{\mathrm p}),\, \Ke(J, C_{\mathrm e})\bigr)
     \,;\; \lstk{L}{s \cons S'}
  \,\Bigr]
  \;\red\;
  \mathcal{C}\Bigl[\,
     \compute(I,J,C_{\mathrm p}, C_{\mathrm e})
     \,;\; \lstk{L}{S'}
  \,\Bigr],
\end{equation}
which fires only when $\near(I,J)=L$ is defined and the head token $s$
matches the signature guard. Firing the rule is \emph{forcing}: it
unwraps the metered thunk by spending exactly one head token. Structural
congruence $\Seq$ (which includes the equations of $K$ and the
location-respecting rearrangements of the bag) consumes no token and
leaves the stack invariant. We import the two structural results
of~\cite{CostNote} that we shall use.

\begin{axiomx}[No leak / wrapping is a sorting invariant; {\cite[Lemma on subject reduction]{CostNote}}]
\label{ax:noleak}
In $\Cost(G)$ every redex is born inside a wrapper, and reduction
preserves wrapping: the only run-time interaction with tokens is
unwrapping. In particular a firing modifies no token other than the head
it spends, and creates no enabling for another redex except through the
contractum it produces.
\end{axiomx}

\begin{axiomx}[Modulus; {\cite[Prop.\ ``stack consumption is the modulus'']{CostNote}}]
\label{ax:modulus}
For a terminating reduction in $\Cost(G)$ the number of stack cells
consumed equals the number of forced redexes, which equals the number of
cuts eliminated --- including those introduced by duplication, each
charged when forced. The consumed prefix of the temporal stack is an
exact operational modulus for the cut elimination.
\end{axiomx}

From the shape of~\eqref{eq:gatedcut} and Property~\ref{ax:noleak} we
extract the single locality fact on which the concurrency analysis turns.
It is a property of the rule, recorded here as the working hypothesis we
verify against~\cite{CostNote} rather than reprove.

\begin{axiomx}[Locality of forcing]
\label{ax:loc}
A firing $\mathcal{C}[\,\rho\,]\red \mathcal{C}[\,\rho'\,]$ that
contracts the cut instance $\rho$ and pops the head of the located stack
at $L$ leaves every subterm of $\mathcal{C}[\,\rho\,]$ disjoint from
$\rho$ and from that head cell unchanged, up to $\Seq$.
\end{axiomx}

\noindent
This is immediate from~\eqref{eq:gatedcut}: the left side names exactly
the contracted cut and the spent head cell, and Property~\ref{ax:noleak}
forbids any ambient side effect, so the rewrite is confined to the named
material.

%% ══════════════════════════════════════════════════════════════════
\section{The pointed domain and firings}
\label{sec:domain}

A causal set needs an origin: the spacetime of a computation is the
spacetime \emph{of a run}, so we point the objects with an initial
configuration.

\begin{definition}[Pointed cost-accounted theory]
\label{def:pointed}
The category $\CostCp$ has as objects pairs $(\Cost(G), t_0)$ with $G$ a
continued interactive GSLT and $t_0$ a closed configuration of
$\Cost(G)$ (a wrapped, signed term carrying located stacks). A morphism
$(\Cost(G), t_0) \to (\Cost(G'), t_0')$ is a $\ciGSLT$-morphism
$f : \Cost(G) \to \Cost(G')$ in the sense of~\cite{CostNote}
(bisimulation-preserving and quote-faithful) with $f(t_0)=t_0'$. We
write $\reach(t_0)$ for the set of configurations reachable from $t_0$
under $\red$, taken modulo $\Seq$ (Property~\ref{ax:loc} makes $\red$
well defined on $\Seq$-classes; see Lemma~\ref{lem:gauge}).
\end{definition}

\begin{definition}[Redex occurrence, firing, footprint]
\label{def:firing}
A \emph{redex occurrence} in a configuration $t$ is the data
$r = (\mathcal{C}, \kappa)$ of a context together with a cut instance
$\kappa = (I,J,C_{\mathrm p},C_{\mathrm e};\, L, s)$ filling its hole, to
which the gated rule~\eqref{eq:gatedcut} applies: $\near(I,J)=L$ and the
head of the stack at $L$ is the matching token $s$. We write
$t \fire{r} t'$ for the corresponding \emph{firing}. Its
\emph{footprint} has two parts, in the sense of Bernstein's conditions
for independence of operations. The \emph{write-set} (the consumed
resources) is the triple
\[
  \foot r \;\defeq\;
  \bigl\{\, \mathrm{intro}_{\mathrm p}(r),\;
            \mathrm{intro}_{\mathrm e}(r),\;
            \mathrm{cell}(r) \,\bigr\},
\]
the two introduction-occurrences $\Kp(I,C_{\mathrm p})$,
$\Ke(J,C_{\mathrm e})$ and the head cell $\mathrm{cell}(r)=(L,s)$ that
the firing spends. The \emph{read-set} $\rfoot r$ is the structural
context on which the rule's applicability depends but which it does
\emph{not} consume: the surfaces $I,J$ inspected by the guard
$\near(I,J)$, the located stack's identity $L$, and --- in calculi where
a step's enabling depends on surrounding structure, such as the boundary
an ambient capability reads or the context in which an \texttt{open}
acts --- that surrounding structure. The contractum and any tokens the
firing pushes are \emph{produced}, part of neither set.
\end{definition}

\begin{lemma}[Forcing is gauge-invariant; quotient $\mathrm{Q1}$]
\label{lem:gauge}
If $t \Seq u$ and $r$ is a redex occurrence of $t$, then the
$\Seq$-corresponding occurrence $r$ of $u$ is a redex occurrence, and
$t \fire{r} t'$, $u \fire{r} u'$ imply $t' \Seq u'$. Hence $\red$,
firings, and footprints are well defined on $\Seq$-classes.
\end{lemma}

\begin{proof}
The gated rule~\eqref{eq:gatedcut} is stated modulo $\Seq$: its premise
($\near(I,J)$ defined, head token matching) and its action (replace the
cut by $\compute$, pop the head) refer only to $\Seq$-classes, since
$\Seq$ leaves both the contact structure up to $K$'s equations and the
located stack invariant. A congruence $t \Seq u$ carries the occurrence
$r$ to a $\Seq$-equal occurrence with identical premise and identical
contractum up to $\Seq$. Confluence of $\Seq$ with $\red$ on the spent
material then gives $t' \Seq u'$.
\end{proof}

\begin{definition}[Independence / concurrency at a configuration]
\label{def:independent}
Two redex occurrences $r_1, r_2$ co-enabled at $t$ are \emph{independent}
(or \emph{concurrent}), written $r_1 \co r_2$, when they satisfy the
\emph{Bernstein conditions}: no write--write overlap and no read--write
overlap, while read--read overlap is permitted:
\[
  \foot r_1 \cap \foot r_2 = \emptyset,
  \qquad
  \foot r_1 \cap \rfoot r_2 = \emptyset,
  \qquad
  \rfoot r_1 \cap \foot r_2 = \emptyset.
\]
That is, neither spends a resource the other spends, and neither spends a
resource the other depends on. The relation $\co$ is irreflexive (a redex
shares its whole write-set with itself) and symmetric.
\end{definition}

A co-enabled pair therefore falls into exactly one of two cases, and the
distinction is the whole point. Either it is independent in the Bernstein
sense above --- and then, by Lemma~\ref{lem:diamond}, the two orders
\emph{converge}: this is genuine concurrency, not a race, since the
outcome is the same whichever fires first. Or the conditions fail, in one
of two ways, each of which is a dependency, not a coincidence to be
smoothed over:
\begin{itemize}[leftmargin=1.6em,itemsep=2pt,topsep=2pt]
  \item \textbf{write--write clash} ($\foot r_1 \cap \foot r_2 \neq
  \emptyset$): both would consume the \emph{same} resource --- the same
  located-stack head, or the same introduction. At one configuration the
  stack has a single head, so two redexes both poised to pop it are
  \emph{competing}; firing either consumes the resource and destroys the
  other redex. This is \emph{conflict}, a race (Section~\ref{sec:events},
  relation $\confl$), and the diamond does not apply to it.
  \item \textbf{read--write clash} ($\foot r_1 \cap \rfoot r_2 \neq
  \emptyset$ or symmetrically): one firing rewrites structural context the
  other reads --- an \texttt{open} dissolving a boundary the other step is
  standing inside. Firing the writer changes whether or how the reader
  fires, so the two are ordered, not interchangeable.
\end{itemize}
Causal order $\le$ between distinct events
(Section~\ref{sec:events}) is a third, separate phenomenon: it relates an
event to a \emph{later} one that consumes or reads what the first
\emph{produced}, across configurations --- not a clash between two redexes
co-enabled at one configuration. The diamond below speaks only to the
first, independent case.

\begin{lemma}[Diamond for independent redexes]
\label{lem:diamond}
If $r_1 \co r_2$ at $t$, then there is a configuration $t_{12}$ with
\[
  t \fire{r_1} t_1 \fire{r_2} t_{12}
  \qquad\text{and}\qquad
  t \fire{r_2} t_2 \fire{r_1} t_{12},
\]
and the write-sets $\foot r_1, \foot r_2$ are unchanged along both
paths. That is, $\red$ is locally confluent \emph{on independent
redexes}. (This is not a confluence claim about $\red$ in general; see
Remark~\ref{rem:nonconfluent}.)
\end{lemma}

\begin{proof}
By Property~\ref{ax:loc}, firing $r_1$ alters only the material in
$\foot r_1$ and its contractum, leaving everything disjoint from
$\foot r_1$ unchanged up to $\Seq$. Since $\foot r_2 \cap \foot r_1 =
\emptyset$ (no write--write overlap), both introductions of $r_2$ and the
head cell of $r_2$ persist in $t_1$ exactly as in $t$. Since $\foot r_1
\cap \rfoot r_2 = \emptyset$ (no read--write overlap), $r_1$ alters
nothing in $r_2$'s read-set: the surfaces $I_2, J_2$ the guard
$\near(I_2,J_2)=L_2$ inspects, the stack identity $L_2$, and any
surrounding structure $r_2$ depends on are untouched, so $r_2$'s guard
still holds. The head token of the stack at $L_2$ is untouched because
$r_1$ spends a different cell (independent co-enabled redexes on a common
stack would both target its single head, violating
$\foot r_1 \cap \foot r_2 = \emptyset$). Hence $r_2$ is enabled at $t_1$,
with the same write-set. Symmetrically $r_1$ is enabled at $t_2$.

For confluence, observe that the two firings rewrite disjoint subterms
of $t$ (their write-sets are disjoint, their reads are not written by the
other, and by Property~\ref{ax:loc} their effects are confined to those
write-sets plus fresh contracta). Rewriting disjoint, non-overlapping
subterms commutes: the result of performing both, in either order, is the
term obtained from $t$ by simultaneously replacing each cut instance by
its contractum and each spent head cell $s_i \cons S_i'$ by its tail
$S_i'$. The two contracta are independent of each other (each is a
function only of its own footprint), and the two stack pops act on
distinct cells, so the outcomes coincide up to $\Seq$. This common
outcome is $t_{12}$.
\end{proof}

\begin{remark}[The calculus is not confluent; all nondeterminism is conflict]
\label{rem:nonconfluent}
Lemma~\ref{lem:diamond} must not be misread as asserting that $\Cost(G)$,
or rho, $\pi$, the ambient calculus, or a general ciGSLT, is confluent.
None of them is, and the construction neither claims nor needs global
confluence. The diamond is restricted to \emph{independent} redexes
(Definition~\ref{def:independent}), and that restriction is exactly the
definition of concurrency, not a determinism hypothesis. A co-enabled
pair that is \emph{not} independent --- a write--write clash --- is a
genuine race, and the diamond is silent about it: the two firings consume
the same resource, each destroying the other, and the run forks. All of
the calculus's nondeterminism is carried by the conflict relation
$\confl$ of the event structure (Section~\ref{sec:events}), never by
$\co$; trace equivalence (Definition~\ref{def:runs}) permutes only
independent steps and so never identifies two resolutions of a race ---
they are distinct events, distinct maximal configurations, distinct
causal sets, distinct spacetimes. Choosing among them is the content of
Section~\ref{sec:conflict}.

The canonical races are write--write clashes on a shared consumed
resource. Two sodium atoms competing for one chlorine: the Cl
introduction lies in both write-sets, $\foot r_{\mathrm{Na}_1} \cap
\foot r_{\mathrm{Na}_2} \ni \mathrm{Cl} \neq \emptyset$, so the pair is
\emph{not} independent; firing either consumes the Cl and annihilates the
other redex, yielding two conflicting events $\mathrm{Na}_1\text{--}
\mathrm{Cl}$ and $\mathrm{Na}_2\text{--}\mathrm{Cl}$ and two
configurations. Two sperm and one egg, or two cores racing to write one
cache line or take one lock, have the identical shape: a shared
write-set, hence $\confl$. By contrast two writes to \emph{different}
locations have disjoint write-sets, genuinely commute, and are not a race
at all --- the only case the diamond touches. So the diamond is precisely
the criterion that tells races from non-races: a race is a co-enabled
pair for which the diamond fails.

There is a deeper reason the split lands here. Spacelike-separated events
must have an order-independent joint outcome, or two observers who
legitimately disagree about which fired first would disagree about who
won, breaking frame-independence (the remark after
Lemma~\ref{lem:histconf}). Spacelike pairs are therefore \emph{forced} to
commute --- the diamond on $\co$ is the discrete shadow of Lorentz
invariance of outcomes at spacelike separation --- so a race cannot be
spacelike. It must be a conflict, and conflict is local: it lives at a
single shared resource both competitors can reach (the Cl, the egg, the
cache line), a locus of decision in causal contact with both. This is the
sense in which races sit ``in the same light cone'': nondeterminism
requires causal contact, and causal contact through a shared consumed
resource is $\confl$, not $\co$. The construction puts races exactly where
relativity says they belong.

Finally, a representational caveat. Prime event structures suffice for
all three races above because each individuates its winner: the events
$\mathrm{Na}_1\text{--}\mathrm{Cl}$ and $\mathrm{Na}_2\text{--}
\mathrm{Cl}$ are distinct, each with its own causal past. What prime
event structures do \emph{not} capture is genuine \emph{OR-causality}
--- a single event with alternative, mutually sufficient enabling sets,
as when one abstracts ``the egg is fertilised'' away from \emph{which}
sperm did it. That requires general (non-prime) event
structures~\cite{WinskelNielsen}. We work with prime event structures
throughout because the cost discipline individuates each forcing by the
cell it spends, so the winner is always named.
\end{remark}

\begin{remark}[Why disjoint footprints, not disjoint occurrences]
The footprint, not the syntactic position, is the right notion of
independence under a non-free $K$. When $K$ is associative--commutative
(the rho case), two introductions floating in the same bag are
positionally adjacent to every stack, so positional disjointness is
vacuous; but they may still spend distinct located cells and hold
distinct surfaces, hence be genuinely concurrent. Conversely under a free
$K$ ($\lambda$) the footprint disjointness reduces to disjointness of the
rigid contexts. The footprint reading is the one that survives the whole
migration spectrum of~\cite{CostNote}, because it is keyed to what is
\emph{consumed}, which is what the cost construction was built to track.
\end{remark}

%% ══════════════════════════════════════════════════════════════════
\section{The event structure of a computation}
\label{sec:events}

We now pass from firings-in-context to \emph{events}: firings identified
across the gauge $\Seq$ (Q1) and across the reordering of concurrent
steps (Q2). The standard machinery is that of asynchronous transition
systems, Mazurkiewicz traces, and prime event
structures~\cite{NPW,WinskelES,WinskelNielsen,Mazurkiewicz,Bednarczyk};
we construct the event structure directly from runs and prove its axioms,
invoking the external theory only for the conceptual packaging.

\begin{definition}[Runs, trace equivalence]
\label{def:runs}
A \emph{run} from $t_0$ is a finite sequence
$\rho = \bigl(t_0 \fire{r_1} t_1 \fire{r_2} \cdots \fire{r_n} t_n\bigr)$
of firings (each $t_i \in \reach(t_0)$, taken modulo $\Seq$). Its length
is $|\rho|=n$. \emph{Trace equivalence} $\sim$ is the least equivalence
on runs from $t_0$ closed under the transposition of adjacent
independent steps: if $r_i \co r_{i+1}$ at $t_{i-1}$, then by
Lemma~\ref{lem:diamond}
\[
  \cdots t_{i-1} \fire{r_i} t_i \fire{r_{i+1}} t_{i+1} \cdots
  \;\sim\;
  \cdots t_{i-1} \fire{r_{i+1}} t_i' \fire{r_i} t_{i+1} \cdots .
\]
A \emph{history} is a $\sim$-class of runs; $\Hist(t_0)$ is the set of
histories, partially ordered by the trace-prefix order
$h \le h'$ iff some representative of $h$ is a prefix of some
representative of $h'$.
\end{definition}

\begin{lemma}[Length is a trace invariant]
\label{lem:length}
If $\rho \sim \rho'$ then $|\rho| = |\rho'|$. Hence $|h| \defeq |\rho|$
for any $\rho \in h$ is well defined.
\end{lemma}

\begin{proof}
A generating transposition permutes two steps, changing neither the
number of steps nor the multiset of firings performed. The general case
follows by induction on the number of transpositions relating $\rho$ to
$\rho'$.
\end{proof}

\begin{definition}[Events]
\label{def:events}
A history $p \in \Hist(t_0)$ is \emph{prime} if it has a greatest proper
sub-history: a unique $p^- \le p$ with $p^- \neq p$ such that
$q \le p,\ q \neq p \Rightarrow q \le p^-$. The set of events is
$\Evt = \Evt(t_0) \defeq \{\,p \in \Hist(t_0) : p \text{ prime}\,\}$,
ordered by restriction of $\le$, with the conflict relation
\[
  e \confl e'
  \quad\defeq\quad
  e, e' \text{ have no common upper bound in } \Hist(t_0).
\]
We write $\Ev(t_0) = (\Evt, \le, \confl)$.
\end{definition}

Intuitively a prime history is a run-modulo-reordering whose final step
cannot be deferred: it is one firing together with exactly the causal
past it requires. The event \emph{is} that minimal history; $e^-$ removes
its final firing.

\begin{lemma}[Prime-event-structure axioms; quotient $\mathrm{Q2}$]
\label{lem:pes}
$\Ev(t_0)$ is a prime event structure:
\begin{enumerate}[label=\textup{(\roman*)},leftmargin=2.2em,itemsep=1pt]
  \item \textbf{Finite causes.} For every $e \in \Evt$ the down-set
  $\dpast e = \{e' \in \Evt : e' \le e\}$ is finite, with
  $|\dpast e| = |e|$.
  \item \textbf{Conflict is irreflexive and symmetric.}
  \item \textbf{Conflict is hereditary:}
  $e \confl e' \le e'' \Rightarrow e \confl e''$.
\end{enumerate}
\end{lemma}

\begin{proof}
\emph{(i)} Each prime $e' \le e$ is determined by a firing occurring in
the history $e$, namely the final firing of $e'$; distinct primes below
$e$ end in distinct firing-occurrences of $e$, and conversely each of the
$|e|$ firing-occurrences in $e$ determines a unique prime below $e$ (its
minimal causal closure). Hence $e' \mapsto (\text{final firing})$ is a
bijection $\dpast e \to \{\text{firings of } e\}$, so $|\dpast e| = |e|$,
which is finite because runs are finite.

\emph{(ii)} Symmetry is immediate from the symmetric phrasing. For
irreflexivity, $e \le e$, so $e$ is a common upper bound of $e$ with
itself; thus $\lnot(e \confl e)$.

\emph{(iii)} Suppose $e \confl e'$ and $e' \le e''$, but for
contradiction $e, e''$ have a common upper bound $u$. Then $e \le u$ and
$e' \le e'' \le u$, so $u$ bounds both $e$ and $e'$, contradicting
$e \confl e'$. Hence $e \confl e''$.
\end{proof}

\begin{definition}[Configurations]
\label{def:config}
For a history $h \in \Hist(t_0)$ put
$\lfloor h\rfloor \defeq \{\, e \in \Evt : e \le h \,\}$. A subset
$x \subseteq \Evt$ is a \emph{configuration} when it is
\emph{down-closed} ($e' \le e \in x \Rightarrow e' \in x$) and
\emph{conflict-free} ($e,e' \in x \Rightarrow \lnot(e \confl e')$). The
set of configurations is $\Conf(\Ev(t_0))$, ordered by inclusion.
\end{definition}

\begin{lemma}[Histories are configurations]
\label{lem:histconf}
For each $h \in \Hist(t_0)$, $\lfloor h\rfloor$ is a configuration with
$|\lfloor h\rfloor| = |h|$. The assignment $h \mapsto \lfloor h\rfloor$
is an order-isomorphism $\Hist(t_0) \cong \Conf(\Ev(t_0))$.
\end{lemma}

\begin{proof}
Down-closure: if $e' \le e \le h$ then $e' \le h$, so $e' \in
\lfloor h\rfloor$. Conflict-freeness: any $e, e' \le h$ have the common
upper bound $h$, so $\lnot(e \confl e')$. The cardinality
$|\lfloor h\rfloor| = |h|$ follows as in Lemma~\ref{lem:pes}(i): the
primes below $h$ are in bijection with the firing-occurrences of $h$,
which number $|h|$ by Lemma~\ref{lem:length}. That $h \mapsto
\lfloor h\rfloor$ is a monotone bijection with monotone inverse is the
prime-algebraicity of the trace-history domain; this is the standard
representation of Mazurkiewicz-trace languages and asynchronous
transition systems by prime event structures and prime algebraic
domains~\cite{NPW,WinskelES,WinskelNielsen,Mazurkiewicz}, whose
hypotheses (the diamond property and finiteness, here
Lemmas~\ref{lem:diamond} and~\ref{lem:length}) we have verified for
$\Cost(G)$.
\end{proof}

\begin{remark}[Frame-independence]
Lemma~\ref{lem:histconf} is the computational content of relativistic
frame-independence. A run is a choice of total order on its firings --- a
choice of observer's clock. Trace equivalence quotients runs by the
reordering of concurrent (spacelike) steps, exactly as different inertial
observers disagree on the time-order of spacelike-separated events. The
invariant is the partial order itself: $\Ev(t_0)$ is the quotient of all
linearisations by permutation of the independent ones, and a
configuration is the frame-free history that all observers agree on.
\end{remark}

%% ══════════════════════════════════════════════════════════════════
\section{Configurations are causal sets}
\label{sec:causet}

\begin{definition}[Causal set; {\cite{BLMS}}]
\label{def:causet}
A \emph{causal set} is a partial order $(C, \preceq)$ that is
\emph{locally finite}: for all $x \preceq y$ the order interval
\[
  [x,y] \defeq \{\, z \in C : x \preceq z \preceq y \,\}
\]
(the discrete \emph{causal diamond}) is finite. A \emph{measured causal
set} is a causal set with the counting measure $\meas(A) = |A|$.
\end{definition}

\begin{theorem}[Every configuration is a causal set]
\label{thm:causet}
For each $x \in \Conf(\Ev(t_0))$, the pair $(x, \le\restrict x)$ is a
causal set under the inherited order, with $\meas = |\cdot|$. In
particular each history $h$ yields a measured causal set
$\lfloor h\rfloor$, and the maximal configurations --- the complete runs
modulo reordering --- are the inextendible discrete spacetimes of the
computation.
\end{theorem}

\begin{proof}
The inherited relation is a partial order. For local finiteness, take
$e \le e'$ in $x$; then $[e,e'] \subseteq \dpast{e'}$, which is finite by
Lemma~\ref{lem:pes}(i). Thus $[e,e']$ is finite. (Note this holds even
when $x$ is infinite, i.e.\ when $t_0$ admits non-terminating runs: the
causal set may be infinite, but every causal diamond in it is finite,
which is exactly local finiteness.)
\end{proof}

The theorem is the structural half of \emph{Order $+$ Number $=$
Geometry}: the order is in hand, and it is of the right kind. The next
section supplies the number.

%% ══════════════════════════════════════════════════════════════════
\section{Number is volume: spent phlogiston is spacetime volume}
\label{sec:numbervolume}

In causal-set theory the continuum spacetime volume of a region is
recovered, in the manifoldlike limit, as the number of causet elements
it contains: volume is cardinality~\cite{BLMS,SorkinValdivia}. We show
that on $\Ev(t_0)$ this cardinality is the spent phlogiston, so that the
modulus of~\cite{CostNote} \emph{is} the volume measure.

\begin{notation}
For an event $e$, let $\spent(e)$ denote the phlogiston spent to reach
$e$: the number of stack cells popped along any run-prefix whose history
is $e$. By Lemma~\ref{lem:length} this is independent of the chosen
linearisation.
\end{notation}

\begin{theorem}[Number $=$ Volume $=$ phlogiston]
\label{thm:numbervolume}
For every $e \in \Evt(t_0)$,
\[
  \meas(\dpast e) \;=\; |\dpast e| \;=\; |e| \;=\; \spent(e).
\]
Consequently, for $e \le e'$ the volume of the causal diamond is the
phlogiston spent strictly within it,
\[
  \meas([e,e']) \;=\; |[e,e']| \;=\; \spent(e') - \spent(e) + 1,
\]
counting both endpoints; and the volume of the causal past of any region
is the total phlogiston its computation has burned.
\end{theorem}

\begin{proof}
The first two equalities are Lemma~\ref{lem:pes}(i). For
$|e| = \spent(e)$: a run-prefix realising the history $e$ performs
exactly $|e|$ firings (Lemma~\ref{lem:length}), and by the modulus
Property~\ref{ax:modulus} each firing is a forcing that pops exactly one
stack cell, the number of cells popped equalling the number of forced
redexes. Hence the cells popped number $|e|$, i.e.\ $\spent(e)=|e|$. The
diamond formula follows since $[e,e'] = \dpast{e'} \setminus
\dpast{e}^{\,\circ}$ where $\dpast{e}^{\,\circ} = \dpast e \setminus
\{e\}$ is the strict past, so $|[e,e']| = |\dpast{e'}| - |\dpast e| + 1
= \spent(e') - \spent(e) + 1$.
\end{proof}

\begin{remark}[Phlogiston as a volume, made literal]
The closing correspondence of~\cite{CostNote} read phlogiston as a
potential on the transition graph. Theorem~\ref{thm:numbervolume} gives
it the complementary geometric reading demanded by causal-set theory:
spent phlogiston is the four-volume swept out by the computation's causal
past. The temporal stack is therefore not merely a clock but the
\emph{volume measure} of the emergent spacetime --- and because the
metering is lazy (charged at forcing, not at exposure), the measure is
accumulated exactly in step with the spacetime's own growth. Burning
phlo is making spacetime.
\end{remark}

%% ══════════════════════════════════════════════════════════════════
\section{Proper time and the Lorentzian signature}
\label{sec:lorentz}

A causal set is Lorentzian, not Riemannian: its order is the causal
(lightcone) order, and the geometric magnitude along a timelike pair is
proper time, characterised in the continuum by the \emph{reverse}
triangle inequality (timelike geodesics maximise proper time). In
causal-set theory the discrete proper time between two related elements
is the length of the longest chain between them~\cite{SorkinValdivia}.
We adopt this and prove the inequality holds here; we then show that
located authority is exactly what produces a nondegenerate signature.

\begin{definition}[Chains, antichains, proper time]
\label{def:chain}
A \emph{chain} in a configuration $x$ is a set of pairwise
$\le$-comparable events (a discrete timelike worldline); an
\emph{antichain} is a set of pairwise incomparable events (a discrete
spacelike slice). For $e \le e'$ the \emph{proper time} is
\[
  \propertime(e, e') \;\defeq\;
  \max\{\, k : e = e_0 < e_1 < \cdots < e_k = e' \,\},
\]
the length of the longest chain in $[e,e']$ (finite by local finiteness,
Theorem~\ref{thm:causet}).
\end{definition}

\begin{proposition}[Lorentzian reverse triangle inequality]
\label{prop:revtriangle}
For $e \le e' \le e''$ in any configuration,
\[
  \propertime(e, e'') \;\ge\; \propertime(e, e') + \propertime(e', e'').
\]
\end{proposition}

\begin{proof}
Let $C_1$ be a longest chain from $e$ to $e'$ (length
$\propertime(e,e')$) and $C_2$ a longest chain from $e'$ to $e''$
(length $\propertime(e',e'')$). Their concatenation at the shared element
$e'$ is a chain from $e$ to $e''$ of length
$\propertime(e,e') + \propertime(e',e'')$. The longest such chain is at
least this long, giving the inequality.
\end{proof}

This is the signature of time: unlike a metric, $\propertime$ adds
\emph{against} subdivision, so inserting an intermediate event can only
lengthen proper time, never shorten it --- the discrete analogue of the
twin paradox. The spatial direction, by contrast, is carried by
antichains, on which $\propertime$ is undefined (incomparable elements
have no chain between them). The question is whether a configuration has
\emph{both}. The answer is governed, as everywhere in~\cite{CostNote}, by
the equational theory of $K$ through the located stack.

\begin{proposition}[Signature from location]
\label{prop:signature}
Fix a configuration $x \in \Conf(\Ev(t_0))$.
\begin{enumerate}[label=\textup{(\alph*)},leftmargin=2.2em,itemsep=2pt]
  \item \emph{(Time from $\cons$.)} The firings in $x$ that draw on a
  single located stack $\lstk{L}{-}$ form a chain: they are pairwise
  $\le$-comparable, ordered by the depth in $L$ at which each pops.
  \item \emph{(Space from $K$.)} If $e_1, e_2 \in x$ draw on distinct
  located stacks and hold surfaces $I_1, I_2$ with $\near(I_1,I_2)$ and
  $\near(I_2,I_1)$ undefined, and neither produces an introduction the
  other consumes, then $e_1, e_2$ are concurrent: they form an antichain.
  \item \emph{(Signature.)} If $x$ contains a located stack drawn on at
  least twice \emph{and} two such non-near surfaces drawn concurrently,
  then $x$ is neither a chain nor an antichain: its ordering fraction
  (Definition~\ref{def:dimension}) lies strictly between $0$ and $1$, the
  combinatorial mark of a $(1{+}n)$-dimensional causal structure.
\end{enumerate}
\end{proposition}

\begin{proof}
\emph{(a)} Two firings spending head cells of the same stack
$\lstk{L}{-}$ cannot be independent, since the head cell is a single
shared resource: at most one of them is enabled to pop the current head
at any configuration (Definition~\ref{def:independent} excludes a common
cell from $\co$). Concretely, if the stack is
$s_1 \cons s_2 \cons \cdots$, then $s_2$ becomes the head only after a
firing pops $s_1$; a firing popping $s_2$ therefore has, in its causal
past, a firing popping $s_1$, so the former $\le$ relates the latter. By
induction the draws on $L$ are linearly ordered by stack depth, hence a
chain.

\emph{(b)} The hypotheses say precisely $\foot{e_1} \cap \foot{e_2} =
\emptyset$: distinct stacks give distinct cells, the non-nearness of the
surfaces means neither's surface is the other's (no shared introduction
through surface matching), and the last clause excludes a
producer/consumer link. Thus $e_1 \co e_2$. Independent events are
incomparable in $\le$ (neither lies in the other's causal past, since
neither's footprint is produced by the other) and compatible (the
diamond, Lemma~\ref{lem:diamond}, gives a common upper bound), so they
are concurrent and lie in a common antichain.

\emph{(c)} By (a) the repeated draws on the stack give a chain of
length $\ge 2$, contributing at least one related pair, so the ordering
fraction $r(x) > 0$ and $x$ is not an antichain. By (b) the two
concurrent non-near draws give an incomparable pair, so $r(x) < 1$ and
$x$ is not a chain. The intermediate value $0 < r(x) < 1$ is, via the
Myrheim--Meyer estimator (Definition~\ref{def:dimension},
Proposition~\ref{prop:dimbounds}), a finite dimension $> 1$.
\end{proof}

\begin{remark}[The welding, restated]
Proposition~\ref{prop:signature} is the formal content of the entangling
of space and time announced in~\cite{CostNote}. Without location, a
configuration would be a bare product: the temporal stacks would order
nothing across surfaces (no shared cells to chain distinct surfaces) and
$K$ would spread surfaces with no temporal coupling, giving an antichain
of independent timelines --- a space and a time that vary independently,
i.e.\ a foliated $\Sigma \times \mathbb{R}$ with a global ``now.''
Location is what binds a temporal chain (a) \emph{to} a spatial surface,
so that a single located stack is a worldline pinned at a place, and the
family of located stacks across non-near surfaces is the spacelike
spread of worldlines. The mixing of (a) and (b) destroys the global
foliation --- there is no slice cutting every worldline at the same
stack-depth, because depths on distinct stacks are incomparable --- and
that is what makes the structure a spacetime rather than a space times a
time.
\end{remark}

\begin{definition}[Ordering fraction; Myrheim--Meyer dimension]
\label{def:dimension}
For a finite configuration or finite causal diamond $A$ with $N = |A|
\ge 2$, let $R(A)$ be the number of \emph{related} (comparable, ordered)
pairs. The \emph{ordering fraction} is
\[
  r(A) \;\defeq\; \frac{R(A)}{\binom{N}{2}}
              \;=\; \frac{2\,R(A)}{N(N-1)} \;\in\; [0,1].
\]
The \emph{Myrheim--Meyer dimension} $d_{\mathrm{MM}}(A)$ is the unique
real $d \ge 1$ solving the sprinkling relation~\cite{Myrheim,MeyerThesis}
\begin{equation}
\label{eq:mm}
  r \;=\; \frac{\Gamma(d+1)\,\Gamma\!\bigl(\tfrac d2\bigr)}
               {4\,\Gamma\!\bigl(\tfrac{3d}{2}\bigr)},
\end{equation}
the expected ordering fraction of a Poisson sprinkling into a causal
diamond of $d$-dimensional Minkowski space.
\end{definition}

\begin{proposition}[Extremal bounds]
\label{prop:dimbounds}
The right-hand side of~\eqref{eq:mm} is strictly decreasing in $d$ on
$[1,\infty)$, with value $1$ at $d=1$ and limit $0$ as $d\to\infty$.
Hence $r(A)=1$ iff $A$ is a chain, forcing $d_{\mathrm{MM}}=1$ (pure
time, no space); $r(A)=0$ iff $A$ is an antichain, the degenerate
$d_{\mathrm{MM}}=\infty$ limit (pure space, no time); and any $A$
satisfying the hypotheses of Proposition~\ref{prop:signature}(c) has
$0 < r(A) < 1$, hence finite $d_{\mathrm{MM}} > 1$.
\end{proposition}

\begin{proof}
$r(A)=1$ means every pair is comparable, i.e.\ $A$ is a chain; $r(A)=0$
means no pair is comparable, i.e.\ $A$ is an antichain. At $d=1$
the diamond of $1$-dimensional Minkowski space is an interval of the
line, whose every sprinkled pair is timelike-related, giving expected
$r=1$; substitution $d=1$ into~\eqref{eq:mm} gives $\Gamma(2)\Gamma(1/2)
/\bigl(4\,\Gamma(3/2)\bigr) = 1\cdot\sqrt\pi/(4\cdot\tfrac{\sqrt\pi}{2})
= 1$. Monotonicity and the $d\to\infty$ limit are the standard
properties of~\eqref{eq:mm}~\cite{MeyerThesis}; we do not reproduce the
asymptotic analysis of the Gamma-quotient here. The final clause is
Proposition~\ref{prop:signature}(c).
\end{proof}

%% ══════════════════════════════════════════════════════════════════
\section{The functor and the spacetime invariant}
\label{sec:functor}

We now assemble the map and identify the equivalence it inverts ---
which is \emph{not} ordinary bisimulation, but its concurrency-aware
refinement.

\begin{definition}[The category of measured causal sets]
\label{def:meascaus}
$\MeasCaus$ has as objects measured causal sets $(C,\preceq,\meas)$ with
$\meas$ the counting measure, and as morphisms the \emph{rigid
measure-preserving causal maps}: injections $g : C \to C'$ that are
order-embeddings ($x \preceq y \iff g(x) \preceq' g(y)$) whose image is
down-closed and which preserve covers ($x \lc y \iff g(x) \lc' g(y)$).
Down-closed image and order-embedding force $\meas([x,y]) =
\meas'([g(x),g(y)])$, so $g$ is automatically measure-preserving on
intervals. Isomorphisms are the order-and-measure isomorphisms.
\end{definition}

The object map sends a pointed system to its event structure, regarded
as the domain of its measured configurations.

\begin{definition}[The map on objects]
\label{def:objmap}
$\Map(\Cost(G), t_0) \defeq \Ev(t_0) = (\Evt, \le, \confl)$ with counting
measure, equivalently its domain of measured configurations
$\Conf(\Ev(t_0))$. A single discrete spacetime is a maximal
configuration $(x, \le\restrict x, |\cdot|) \in \MeasCaus$
(Theorem~\ref{thm:causet}).
\end{definition}

To make $\Map$ a functor we must say which $\CostCp$-morphisms it sends
to causal maps, and here a subtlety of concurrency theory is decisive.
Ordinary (interleaving) bisimulation equates systems with the same tree
of possible step-sequences; it is blind to whether two enabled steps are
concurrent or in conflict, because both present as a branch. But
concurrency is exactly the space in our spacetime: an antichain and a
chain can have the same interleaving tree. Hence the spacetime cannot be
an invariant of interleaving bisimulation. The right notion is the one
that additionally matches the independence relation along the
bisimulation: \emph{hereditary history-preserving bisimulation} (HHPB),
which van Glabbeek and Goltz~\cite{vGG} identify as precisely the
behavioural equivalence corresponding to isomorphism of event
structures.

\begin{definition}[Quote-faithful HHP-morphism]
\label{def:hhpb}
A $\CostCp$-morphism $f : (\Cost(G), t_0) \to (\Cost(G'), t_0')$ is
\emph{footprint-preserving} when it is quote-faithful (preserves
signatures and surfaces, hence the definedness of $\near$ and the
read-set $\rfoot r$) and sends each firing $r$ to a firing $f(r)$ with
$\foot{f(r)} = f(\foot r)$ and $\rfoot{f(r)} = f(\rfoot r)$,
bijectively on both consumed resources (the write-set) and depended-on
context (the read-set). It is an \emph{HHP-morphism} when, in addition,
the induced relation on runs is a history-preserving bisimulation: it
respects enabling and independence, and back-and-forth along
already-matched causal pasts. Write $\CostCp^{\,\mathrm{hhp}}$ for the
subcategory of $\CostCp$ on these morphisms.
\end{definition}

\begin{lemma}[HHP-morphisms preserve the causal structure]
\label{lem:preserve}
A footprint-preserving HHP-morphism $f$ induces a map
$\Ev(f) : \Ev(t_0) \to \Ev(t_0')$ on events that preserves and reflects
$\le$, preserves and reflects $\confl$, and preserves the counting
measure: $\meas(\dpast{e}) = \meas'(\dpast{\Ev(f)(e)})$.
\end{lemma}

\begin{proof}
Because $f$ is footprint-preserving, it carries co-enabled redex
occurrences to co-enabled occurrences with $\foot{f(r)} = f(\foot r)$ and
$\rfoot{f(r)} = f(\rfoot r)$, and bijectivity on both sets carries each
Bernstein clause to its image: write--write, read--write, and
write--read overlaps are preserved and reflected. Hence
$r_1 \co r_2 \iff f(r_1) \co f(r_2)$, so $f$ preserves and reflects
independence (and likewise the failures of independence, i.e.\ races).
Preserving enabling and independence, $f$
sends runs to runs and trace-equivalent runs to trace-equivalent runs
(it commutes with the generating transpositions), so it descends to
$\Hist(t_0) \to \Hist(t_0')$, monotone for the prefix order, carrying
primes to primes; this is $\Ev(f)$. Reflection of $\le$ and the
back-condition of an HHP-bisimulation give injectivity and order-reflection.
Conflict is preserved and reflected because $f$ preserves and reflects
the existence of common upper bounds (it is a bijection on the matched
configuration lattices, by the hereditary back-and-forth). Finally, by
Lemma~\ref{lem:pes}(i) and Lemma~\ref{lem:length}, $\meas(\dpast e) =
|e|$ counts firings, and $f$ is a bijection on the firings of $e$ onto
those of $f(e)$, so the measure is preserved.
\end{proof}

\begin{theorem}[Functoriality]
\label{thm:functor}
$\Map : \CostCp^{\,\mathrm{hhp}} \to \MeasCaus$, sending
$(\Cost(G),t_0) \mapsto \Ev(t_0)$ and $f \mapsto \Ev(f)$, is a functor.
\end{theorem}

\begin{proof}
$\Ev(f)$ is a rigid measure-preserving causal map by
Lemma~\ref{lem:preserve} (order-embedding with down-closed image from
preservation/reflection of $\le$ and the back-condition; cover
preservation from preservation of the immediate-enabling step).
Identities: $\Map(\id) = \Ev(\id)$ fixes every run, hence every history
and every event, so $\Ev(\id) = \id$. Composition: for composable $f, g$,
the induced action on runs satisfies $\Ev(g \circ f) = \Ev(g) \circ
\Ev(f)$ since both are computed stepwise on firings and
$(g\circ f)(\foot r) = g(f(\foot r))$. Hence $\Map$ preserves identities
and composition.
\end{proof}

\begin{corollary}[The spacetime is an HHPB-invariant; quotient $\mathrm{Q3}$]
\label{cor:invariance}
If $(\Cost(G), t_0)$ and $(\Cost(G'), t_0')$ are related by a
quote-faithful hereditary history-preserving bisimulation, then
$\Ev(t_0) \cong \Ev(t_0')$ as measured causal sets, and hence their
domains of measured configurations are isomorphic. In particular the
discrete spacetime of a cost-accounted computation is well defined
independently of the presentation, up to HHPB.
\end{corollary}

\begin{proof}
A quote-faithful HHPB is an HHP-morphism with an HHP-inverse; by
Theorem~\ref{thm:functor} both induce causal maps, and functoriality
makes the induced maps mutually inverse isomorphisms in $\MeasCaus$, by
van Glabbeek--Goltz's correspondence between HHPB and event-structure
isomorphism~\cite{vGG}.
\end{proof}

\begin{remark}[Interleaving bisimulation is the wrong invariant, and why it matters]
It is worth stating the negative result plainly, since it corrects a
natural expectation. The $\ciGSLT$-morphisms of~\cite{CostNote} are
defined up to bisimulation; but \emph{interleaving} bisimulation
identifies a chain of two firings with a confluent diamond of two
concurrent firings whenever their interleaving trees coincide, and these
have different causal sets --- a $2$-chain (timelike) versus a
$2$-antichain (spacelike). Were the spacetime an interleaving-bisimulation
invariant, time and space would be indistinguishable. They are
distinguished exactly by the independence relation, which the footprint
makes available and which HHPB preserves. The located stack is therefore
doing double duty: operationally it prevents ambient authority
(\cite{CostNote}, Section~11), and semantically it is the carrier of the
independence relation that makes the spacetime invariant well defined. The
ocap discipline and the Lorentzian structure are the same structure seen
twice.
\end{remark}

%% ══════════════════════════════════════════════════════════════════
\section{The field-equation shape: naturality of the cover}
\label{sec:naturality}

The general-relativity correspondence of~\cite{CostNote} was a two-way
coupling: the information distribution shapes the nearness structure, and
the nearness structure selects the next redistribution. We render the
coupling as a property of the configuration domain --- that one reduction
step is one cover, whose new geometry is a \emph{local functional} of the
information consumed.

\begin{proposition}[One step is one cover; the increment of geometry is local]
\label{prop:naturality}
Let $x \lc x'$ be a cover in $\Conf(\Ev(t_0))$ (so $x' = x \cup \{e\}$
for a single new event $e$, by Lemma~\ref{lem:histconf} and
Lemma~\ref{lem:length}). Then:
\begin{enumerate}[label=\textup{(\roman*)},leftmargin=2.2em,itemsep=1pt]
  \item $x \lc x'$ corresponds to exactly one firing $r$ with $e = $ the
  prime generated by $r$ in the history $x'$;
  \item the new causal links into $e$ --- the events $e'' \lc e$ ---
  are precisely the events of $x$ that produced the introductions and
  exposed the head cell in $\foot r$; in particular the link-set of $e$
  is a bounded functional of $\foot r$ alone (at most two
  producer-events for the introductions and one for the cell);
  \item the firings enabled at $x'$ but not at $x$ are exactly those
  whose footprint is met by the contractum of $r$ (its produced
  introductions and pushed cells).
\end{enumerate}
Hence ``geometry sourced by the distribution'' is \textup{(ii)} --- the
local information consumed determines the new event's lightcone links ---
and ``distribution guided by geometry'' is \textup{(iii)} --- the new
local geometry determines which interactions come next. The map ``perform
one reduction'' $\mapsto$ ``add one measured cover'' is natural in $\red$.
\end{proposition}

\begin{proof}
\emph{(i)} Adding one firing to a history lengthens it by one
(Lemma~\ref{lem:length}); by Lemma~\ref{lem:histconf} the corresponding
configurations differ by one event, the prime whose final firing is $r$.
\emph{(ii)} By definition $e' \lc e$ iff $e' < e$ with nothing strictly
between; $e' < e$ holds iff $e'$ lies in the causal past required to
enable $r$, which by Property~\ref{ax:loc} and
Definition~\ref{def:firing} means $e'$ produced a member of $\foot r$
(an introduction $r$ contracts, or the cell $r$ pops); immediacy selects
the most recent such producer per resource. There are at most three
resources in $\foot r$, bounding the link-set. \emph{(iii)} A firing $q$
is newly enabled at $x'$ iff its footprint, absent or incomplete at $x$,
is completed by what $r$ produced; by Property~\ref{ax:loc} the only
new material is the contractum of $r$, so $q$'s footprint must meet it.
Naturality is the statement that the assignment commutes with composition
of reduction steps, immediate from (i)--(iii) applied stepwise.
\end{proof}

\begin{remark}[Discrete Einstein coupling]
Proposition~\ref{prop:naturality} is the discrete, computational shadow
of $G_{\mu\nu} = 8\pi T_{\mu\nu}$ read as Wheeler's two-way slogan:
\emph{matter tells spacetime how to curve; spacetime tells matter how to
move}. Here the consumed information (the footprint, the local stress) is
$T$; the new event's link-set (the local causal geometry it creates) is
the local increment of the metric; and (ii)--(iii) are the two directions
of the coupling, each \emph{local and bounded}, as a field equation
demands. We do not claim a quantitative field equation --- that would
require the manifoldlike limit and a continuum action, and is the open
problem flagged in~\cite{CostNote}. What is proved is the \emph{shape}:
the geometry-information coupling is a local functional, step-by-step,
exactly as a hyperbolic field equation is.
\end{remark}

%% ══════════════════════════════════════════════════════════════════
\section{The four quotients and the manifold}
\label{sec:quotients}

We can now name precisely the senses in which the target is a
\emph{quotiented} spacetime. The construction performs four
identifications, the first three of which deliver the causal set as an
invariant and the fourth of which delivers the continuum manifold.

\begin{enumerate}[label=\textbf{(Q\arabic*)},leftmargin=3.2em,itemsep=3pt]
\item \textbf{Gauge: structural congruence $\Seq$.} Re-bracketing the
bag, $\alpha$-renaming, unit insertions --- congruence steps consume no
token and leave the stack invariant. They move the description, not the
events. These are the coordinate changes of the discrete spacetime;
firings act on $\Seq$-classes (Lemma~\ref{lem:gauge}).

\item \textbf{Frame: concurrent permutation (trace equivalence).} Runs
differing only by the order of concurrent (spacelike) firings are the
same history (Definition~\ref{def:runs}); the quotient of all
linearisations by permutation of the independent ones is the causal
partial order (Lemma~\ref{lem:histconf}). This is relativistic
frame-independence.

\item \textbf{Presentation: hereditary history-preserving bisimulation.}
Quote-faithful HHPB-equivalent systems have isomorphic measured causal
sets (Corollary~\ref{cor:invariance}); this is the level on which $\Map$
is a functor, and it is strictly finer than interleaving bisimulation,
which would collapse the space.

\item \textbf{Continuum: the Hauptvermutung class.} A manifoldlike causal
set is identified with the family of Lorentzian manifolds into which it
faithfully embeds, unique up to approximate isometry
(Definition~\ref{def:manifold} below).
\end{enumerate}

\begin{definition}[Faithful embedding; manifoldlike; Hauptvermutung quotient; {\cite{BLMS,BombelliNoldus,MajorRideoutSurya}}]
\label{def:manifold}
A causal set $(C,\preceq)$ \emph{faithfully embeds} into a Lorentzian
manifold $(M,g)$ at density $\varrho$ when there is an order-preserving
injection $\iota : C \hookrightarrow M$ (relating $\preceq$ to the causal
order of $g$) such that $C$ is, with high probability, a Poisson
sprinkling of $M$ at density $\varrho$ --- equivalently, the counting
measure approximates $\varrho$ times the volume measure on regions large
compared with $\varrho^{-1}$. $C$ is \emph{manifoldlike} when such an
$(M,g)$ exists. The \emph{Hauptvermutung} of causal-set theory asserts
that the manifold is then unique up to approximate isometry; we define
\[
  \mathrm{Manifold}(C) \;\defeq\; [\,(M,g)\,]_{\approx},
\]
the approximate-isometry class, when $C$ is manifoldlike, and undefined
otherwise. A precise closeness measure realising $\approx$ is Noldus's
Lorentzian Gromov--Hausdorff distance~\cite{BombelliNoldus}, and the
recovery of continuum topology and dimension is treated
in~\cite{MajorRideoutSurya,MeyerThesis,Myrheim}.
\end{definition}

\begin{definition}[The full map]
\label{def:fullmap}
For $(\Cost(G), t_0) \in \CostCp^{\,\mathrm{hhp}}$ and a chosen maximal
configuration $x$ (one consistent history; see
Section~\ref{sec:conflict} on the choice), the spacetime map is the
composite
\[
  (\Cost(G), t_0)
  \;\xmapsto{\;\Map\;}\;
  \Ev(t_0)
  \;\xmapsto{\;\text{select } x\;}\;
  (x, \le\restrict x, |\cdot|)
  \;\xmapsto{\;[\,\cdot\,]_{\approx}\;}\;
  \mathrm{Manifold}(x),
\]
the last arrow defined precisely when $x$ is manifoldlike. The
``(quotiented) spacetime manifold'' of the computation is
$\mathrm{Manifold}(x)$; the robust, always-defined invariant beneath it
is the measured causal set $(x,\le\restrict x,|\cdot|)$, well defined up
to \textup{(Q1)--(Q3)}.
\end{definition}

\begin{remark}[The manifold is a limit, not a guarantee]
We stress the honest reading. The clean theorem is
$\Map : \CostCp^{\,\mathrm{hhp}} \to \MeasCaus$
(Theorem~\ref{thm:functor}): every cost-accounted computation has a
measured causal set, functorially and HHPB-invariantly. The manifold is
the further \emph{idealisation} of Definition~\ref{def:manifold},
available only for manifoldlike configurations and even then only up to
the Hauptvermutung. Most configurations of a generic $\Cost(G)$ will not
be manifoldlike --- they will be ``too ordered'' (long chains, low
dimension, from a single heavily-drawn stack) or ``too disordered''
(wide antichains, from many non-interacting surfaces). The interesting
question, parallel to the central question of causal-set dynamics, is
which $\Cost(G)$ and which $t_0$ produce manifoldlike spacetimes, and
with what effective dimension --- which by Proposition~\ref{prop:signature}
is a question about the equational theory of $K$ and the discipline of
location.
\end{remark}

%% ══════════════════════════════════════════════════════════════════
\section{The conflict seam: one spacetime or a sum over them}
\label{sec:conflict}

Definition~\ref{def:fullmap} selected a single maximal configuration.
The selection is where conflict ($\confl$) --- the third relation of the
event structure --- enters, and it is exactly the seam between the two
correspondences of~\cite{CostNote}.

A maximal \emph{conflict-free} down-closed set of events is one
consistent history: its image is a single measured causal set, a single
classical spacetime. This is the \emph{general-relativity reading}: a
decohered, definite geometry. The branching of $\confl$ --- competition
of two receivers for one located stack (\cite{CostNote}, Section~11) ---
is the choice of which spacetime.

Keeping the conflict, rather than choosing through it, and weighting
events in a semiring $R$, gives the second reading. The semiring ladder
of~\cite{CostNote} climbs from the Boolean semiring (bare
nondeterminism), through the tropical semiring $(\mathbb{R}\cup\{\infty\},
\min,+)$ (cost; the tropical path-weighting refines the proper time of
Section~\ref{sec:lorentz} from longest \emph{chain} to longest
\emph{weighted} chain), to the non-negative reals (probabilistic), to the
complex numbers (\emph{amplitudes}). At $R=\mathbb{C}$ the object is the
whole branching ensemble and the amplitude of an outcome sums over the
conflicting histories reaching it --- the \emph{quantum-field reading},
in the lineage of Girard's geometry of interaction and the categorical
quantum mechanics of Abramsky and Coecke~\cite{GirardGoI,AbramskyCoecke}.

Here the obstruction flagged in~\cite{CostNote} reappears precisely and
must be stated, not hidden. Interference \emph{cancels} contributions to
a shared outcome; but the configuration domain of an event structure is a
\emph{stable} family, and the equivalence we used to make the spacetime
an invariant, HHPB, is a branching-time notion that forbids exactly such
cancellation~\cite{vGG,WinskelNielsen}. The genuinely quantum reading
therefore pulls the equivalence from branching time toward a linear-time,
trace-indexed semantics --- a reformulation, not a relabelling. We
record this as the boundary of the present construction: the functor
$\Map$ of Theorem~\ref{thm:functor} lands the GR reading (one
configuration, one causal set, HHPB-invariant); the QFT reading requires
replacing $\MeasCaus$-valued configurations by an $R$-weighted,
linear-time object on which cancellation is admissible, and is left, as
in~\cite{CostNote}, to a separate development.

%% ══════════════════════════════════════════════════════════════════
\section{Summary of the construction}
\label{sec:summary}

We collect the specification in one place. The map
$\Map : \CostCp^{\,\mathrm{hhp}} \to \MeasCaus$ is defined by:

\begin{center}
\renewcommand{\arraystretch}{1.35}
\begin{tabular}{@{}>{\raggedright\arraybackslash}p{0.30\textwidth}>{\raggedright\arraybackslash}p{0.62\textwidth}@{}}
\toprule
\textbf{Spacetime datum} & \textbf{Cost-accounted origin} \\
\midrule
event & a cut firing (forcing): Def.~\ref{def:firing} \\
causal order $e \le e'$ & enabling: $e'$ consumes what $e$ produced /
exposed: \S\ref{sec:events} \\
conflict $e \confl e'$ (race) & write--write clash: competition for one
consumed resource (stack head, egg, cache line):
Def.~\ref{def:events}, Rem.~\ref{rem:nonconfluent} \\
concurrency (spacelike) & Bernstein independence (no write--write or
read--write overlap) $e \co e'$: Def.~\ref{def:independent} \\
measure / volume & spent phlogiston $= |{\dpast e}|$:
Thm.~\ref{thm:numbervolume} \\
proper time & longest chain; reverse triangle:
Prop.~\ref{prop:revtriangle} \\
timelike worldline & draws on one located stack $\lstk{L}{-}$:
Prop.~\ref{prop:signature}(a) \\
spacelike slice & non-near surfaces drawn concurrently:
Prop.~\ref{prop:signature}(b) \\
$(1{+}n)$ signature & located authority mixing chains and antichains:
Prop.~\ref{prop:signature}(c) \\
field-equation coupling & link-set of new event is a local functional of
$\foot r$: Prop.~\ref{prop:naturality} \\
the manifold & Hauptvermutung class of a manifoldlike configuration:
Def.~\ref{def:manifold} \\
\bottomrule
\end{tabular}
\end{center}

\noindent
and the ``(quotiented)'' resolves into the four identifications
\textup{(Q1)} $\Seq$ gauge, \textup{(Q2)} trace/frame, \textup{(Q3)}
HHPB/presentation, \textup{(Q4)} Hauptvermutung/continuum
(\S\ref{sec:quotients}). The robust output is the measured causal set,
proved to exist functorially for every pointed cost-accounted theory; the
manifold is its continuum idealisation when manifoldlike.

The single structural lever throughout is the one~\cite{CostNote}
isolated: the equational theory of $K$ and its consequence, the located
stack. A free $K$ gives separable, foliated structure (independent
timelines); any equation on $K$ forces the nominal surface and the
located stack, which welds a temporal chain to a spatial surface and so
makes the structure a spacetime. The object-capability discipline and the
Lorentzian signature are, as Section~\ref{sec:functor} shows, two faces
of the same located authority.

%% ══════════════════════════════════════════════════════════════════
\begin{thebibliography}{99}
\addcontentsline{toc}{section}{References}

\bibitem{CostNote}
L.~G.~Meredith,
\emph{Continued Interactive GSLTs and the Cost Endofunctor},
F1R3FLY.io, 2026.
\url{https://github.com/F1R3FLY-io/publications/blob/main/cost-accounting-as-monad/continued-gslt-cost-v2.pdf}

\bibitem{BLMS}
L.~Bombelli, J.~Lee, D.~Meyer, R.~D.~Sorkin,
\emph{Space-time as a causal set},
Physical Review Letters \textbf{59} (1987) 521--524.

\bibitem{SorkinValdivia}
R.~D.~Sorkin,
\emph{Causal sets: Discrete gravity},
in \textit{Lectures on Quantum Gravity} (Valdivia 2002),
A.~Gomberoff and D.~Marolf, eds., Springer, 2005;
arXiv:gr-qc/0309009.

\bibitem{Myrheim}
J.~Myrheim,
\emph{Statistical geometry},
CERN preprint TH-2538, 1978.

\bibitem{MeyerThesis}
D.~A.~Meyer,
\emph{The Dimension of Causal Sets},
Ph.D.\ thesis, Massachusetts Institute of Technology, 1988.

\bibitem{BombelliNoldus}
J.~Noldus,
\emph{A Lorentzian Lipschitz, Gromov--Hausdorff notion of distance},
Classical and Quantum Gravity \textbf{21} (2004) 839--850;
L.~Bombelli and J.~Noldus,
\emph{The moduli space of causal sets},
Classical and Quantum Gravity \textbf{21} (2004) 4429.

\bibitem{MajorRideoutSurya}
S.~Major, D.~Rideout, S.~Surya,
\emph{On recovering continuum topology from a causal set},
Journal of Mathematical Physics \textbf{48} (2007) 032501.

\bibitem{NPW}
M.~Nielsen, G.~Plotkin, G.~Winskel,
\emph{Petri nets, event structures and domains, part I},
Theoretical Computer Science \textbf{13} (1981) 85--108.

\bibitem{WinskelES}
G.~Winskel,
\emph{Event structures},
in \textit{Petri Nets: Applications and Relationships to Other Models of
Concurrency}, Lecture Notes in Computer Science \textbf{255},
Springer, 1987, 325--392.

\bibitem{WinskelNielsen}
G.~Winskel, M.~Nielsen,
\emph{Models for concurrency},
in \textit{Handbook of Logic in Computer Science}, vol.~4,
S.~Abramsky, D.~Gabbay, T.~Maibaum, eds., Oxford University Press, 1995,
1--148.

\bibitem{Mazurkiewicz}
A.~Mazurkiewicz,
\emph{Trace theory},
in \textit{Petri Nets: Applications and Relationships to Other Models of
Concurrency}, Lecture Notes in Computer Science \textbf{255},
Springer, 1987, 279--324.

\bibitem{Bednarczyk}
M.~A.~Bednarczyk,
\emph{Categories of Asynchronous Systems},
Ph.D.\ thesis, University of Sussex, 1987.

\bibitem{vGG}
R.~J.~van Glabbeek, U.~Goltz,
\emph{Refinement of actions and equivalence notions for concurrent
systems},
Acta Informatica \textbf{37} (2001) 229--327.

\bibitem{Levy}
J.-J.~L\'evy,
\emph{R\'eductions correctes et optimales dans le lambda-calcul},
Th\`ese de doctorat d'\'etat, Universit\'e Paris VII, 1978.

\bibitem{GirardGoI}
J.-Y.~Girard,
\emph{Geometry of interaction I: Interpretation of System F},
in \textit{Logic Colloquium '88}, North-Holland, 1989, 221--260.

\bibitem{AbramskyCoecke}
S.~Abramsky, B.~Coecke,
\emph{A categorical semantics of quantum protocols},
in \textit{Proc.\ 19th IEEE Symposium on Logic in Computer Science
(LICS)}, 2004, 415--425.

\end{thebibliography}

\end{document}
